\documentclass[a4paper,landscape,8pt]{article}

% Packages
\usepackage[landscape,margin=0.5cm]{geometry}
\usepackage{multicol}
\usepackage{amsmath,amssymb,amsfonts}
\usepackage{graphicx}
\usepackage{enumitem}
\usepackage{array}
\usepackage{color}
\usepackage{xcolor}
\usepackage{tcolorbox}
\usepackage{tikz}
\usepackage{bm}

% Custom settings
\setlength{\columnsep}{0.5cm}
\setlength{\columnseprule}{0.2pt}
\setlength{\parindent}{0pt}
\setlength{\parskip}{1pt}
\pagestyle{empty}

% Define colors
\definecolor{heading}{RGB}{0,51,102}
\definecolor{subheading}{RGB}{0,90,120}
\definecolor{example}{RGB}{100,0,0}

% Custom environments and commands
\newcommand{\heading}[1]{\noindent\textbf{\textcolor{heading}{\large #1}}\\\rule{\linewidth}{1pt}}
\newcommand{\subheading}[1]{\vspace{1pt}\noindent\textbf{\textcolor{subheading}{#1}}}
\newcommand{\definition}[2][]{\vspace{1pt}\noindent%
  \ifx\relax#1\relax%
    #2%
  \else%
    \textbf{#1 :} #2%
  \fi}
\newcommand{\theorem}[2]{\vspace{1pt}\noindent\textbf{Theorem (#1):} #2}
\newcommand{\example}[1]{\vspace{1pt}\noindent\textit{\textcolor{example}{Example:} #1}}

\newcommand{\gameMatrix}[1]{
  \left(\begin{array}{cc}
    #1
  \end{array}\right)
}

\begin{document}

\begin{multicols}{3}

	% CHAPTER 1: STATIC GAMES OF COMPLETE INFORMATION
	\heading{1. Static Games of Complete Information}

	\subheading{Normal-Form Game Representation}\\
	\definition{}{The normal-form (strategic-form) representation of an $n$-player game specifies the strategy spaces $S_i$ and payoff functions $u_i$, denoted as $G = \{S_1, \ldots, S_n; u_1, \ldots, u_n\}$}

	\subheading{Dominant Strategies}\\
	\definition{In game $G$, strategy $s_i'$ \textit{strictly dominates} strategy $s_i''$ if:
		\begin{equation*}
			u_i(s_i', s_{-i}) > u_i(s_i'', s_{-i}), \quad \forall s_{-i} \in S_{-i}
		\end{equation*}}


	\subheading{Nash Equilibrium}\\
	\definition[Best Response]{The best response for player $i$ to strategies $s_{-i}$ is:
		\begin{equation*}
			R_i(s_{-i}) = \arg\max_{s_i \in S_i} u_i(s_i, s_{-i})
		\end{equation*}}

	\definition[Nash equilibrium]{Strategies $(s_1^*, \ldots, s_n^*)$ form a Nash equilibrium if:
		\begin{equation*}
			s_i^* \in R_i(s_{-i}^*), \quad \forall i = 1, \ldots, n
		\end{equation*}
		Equivalently: $u_i(s_i^*, s_{-i}^*) = \max_{s_i \in S_i} u_i(s_i, s_{-i}^*)$}


	\subheading{Mixed Strategies}\\
	\definition{A mixed strategy for player $i$ is a probability distribution $p_i = (p_{i1}, \ldots, p_{iK})$ over pure strategies, where $\sum_k p_{ik} = 1$ and $p_{ik} \geq 0$}

	\definition[Expected Payoff]{For a 2-player game:
		\begin{align*}
			v_1(p_1, p_2) & = \sum_{j=1}^J\sum_{k=1}^K p_{1j}p_{2k}u_1(s_{1j}, s_{2k}) \\
			v_2(p_1, p_2) & = \sum_{j=1}^J\sum_{k=1}^K p_{1j}p_{2k}u_2(s_{1j}, s_{2k})
		\end{align*}}

	\definition{Mixed strategies $(p_1^*, p_2^*)$ form a Nash equilibrium if:
		\begin{align*}
			v_1(p_1^*, p_2^*) & \geq v_1(p_1, p_2^*) \\
			v_2(p_1^*, p_2^*) & \geq v_2(p_1^*, p_2)
		\end{align*}
		for all probability distributions $p_1$ and $p_2$}

	\theorem{Nash (1950)}{In any finite $n$-player normal-form game, there exists at least one Nash equilibrium, possibly involving mixed strategies.}

	% CONTENT FOR COLUMN 2 WILL GO HERE (Chapter 2, etc.)
	% CHAPTER 2: DYNAMIC GAMES OF COMPLETE INFORMATION
	\heading{2. Dynamic Games of Complete Information}

	\subheading{Dynamic Games}\\
	\definition{A dynamic game of complete and perfect information is a game where:
		\begin{itemize}[leftmargin=10pt,itemsep=0pt]
			\item Players move in sequence
			\item Perfect information: all previous moves are observed before next move is chosen
			\item Complete information: payoffs are common knowledge
		\end{itemize}
		Such games can be represented by a game tree.}

	\subheading{Backwards Induction}
	\definition{
		\begin{itemize}[leftmargin=10pt,itemsep=0pt]
			\item At the second stage, player 2 observes $a_1$ and solves $\max_{a_2 \in A_2} u_2(a_1, a_2)$
			\item The solution is player 2's best response $R_2(a_1)$
			\item Knowing this, player 1 solves $\max_{a_1 \in A_1} u_1(a_1, R_2(a_1))$
			\item The solution $a^*_1$ and $R_2(a^*_1)$ is the backwards-induction outcome
		\end{itemize}}

	\subheading{Extensive-Form Games and Strategies}\\
	\definition{The extensive-form representation of a game specifies:
		\begin{itemize}[leftmargin=10pt,itemsep=0pt]
			\item The players in the game
			\item When each player has the move, what each player can do, and what each player knows when moving
			\item The payoffs received for each combination of moves
		\end{itemize}}

	\definition[Information Set]{A collection of decision nodes satisfying:
		\begin{itemize}[leftmargin=10pt,itemsep=0pt]
			\item The player needs to move at every node in the set
			\item When play reaches a node in the set, the player does not know which node in the set has been reached
			\item The player must have the same actions available at each node in the set
		\end{itemize}}

	\definition[Strategy]{A complete plan of action specifying a feasible action for the player in every contingency where the player might be called on to act.}


	\subheading{Subgame-Perfect Nash Equilibrium}

	\definition[Subgame]{A portion of a game that:
		\begin{itemize}[leftmargin=10pt,itemsep=0pt]
			\item Begins at a decision node that is a singleton information set
			\item Includes all nodes following it in the game tree
			\item Does not cut any information sets
		\end{itemize}}

	\definition{A Nash equilibrium is \textbf{subgame-perfect} if the players' strategies constitute a Nash equilibrium in every subgame.}

	% CONTENT FOR COLUMN 3 WILL GO HERE (Chapter 3, etc.)
	% CHAPTER 3: STATIC GAMES OF INCOMPLETE INFORMATION
	\heading{3. Static Games of Incomplete Information}

	\subheading{Bayesian Games}
	\definition{Games of incomplete information are also called Bayesian games. In a game of incomplete information, at least one player is uncertain about another player's payoff function.}

	\subheading{Static Bayesian Games and Bayesian Nash Equilibrium} \\
	\definition{The normal-form representation of an $n$-player static Bayesian game specifies:
	\begin{itemize}[leftmargin=10pt,itemsep=0pt]
		\item The players' action spaces $A_1, \ldots A_n$,
		\item Their type spaces $T_1, \ldots, T_n$,
		\item Their beliefs $P_1, \ldots, P_n$, and
		\item Their payoff functions $u_1, \ldots, u_n$.
	\end{itemize}
	\begin{itemize}[leftmargin=10pt,itemsep=0pt]
		\item Moreover, Player $i$'s type, $t_i \in T_i$, is privately known by player $i$, and determines player $i$'s payoff function, $u_i(a_1, \ldots, a_n; t_i)$.}
		\item Player $i$'s belief $P_i(t_{-i}|t_i)$ describes $i$'s uncertainty about the $n - 1$ other players' possible types, $t_{-i}$, given $i$'s own type, $t_i$.
	\end{itemize}

	\definition[Strategy]{In the static Bayesian game, a strategy for player $i$ is a function $s_i: T_i \rightarrow A_i$.
		\begin{itemize}[leftmargin=10pt,itemsep=0pt]
			\item For given type $t_i$, $s_i(t_i)$ gives the action.
			\item Player $i$'s strategy space $S_i$ is the set of all functions from $T_i$ into $A_i$.
		\end{itemize}}

	\definition{In the static Bayesian game, the strategies $s^* = (s_1^*, \ldots, s_n^*)$ are a (pure-strategy) \textbf{Bayesian Nash equilibrium} if for each player $i$ and for each of $i$'s type $t_i$ in $T_i$, $s_i^*(t_i)$ solves
		\begin{equation*}
			\max_{a_i \in A_i} E_{t_{-i}}u_i(a_i, s_{-i}^*(t_{-i}); t_i),
		\end{equation*}
		where
		\begin{equation*}
			E_{t_{-i}} u_i(a_i, s_{-i}^*(t_{-i}); t_i) = \sum_{t_{-i} \in T_{-i}} P_i(t_{-i}|t_i)u_i(a_i, s_{-i}^*(t_{-i}); t_i).
		\end{equation*}}

	\theorem{Proposition 1}{Let $(s_1^*, \ldots, s_n^*)$ be strategies in a static Bayesian game. If for any $t_i \in T_i$, $a_i \in A_i$ and $a_{-i} \in A_{-i}$,
		\begin{equation*}
			u_i(s_i^*(t_i), a_{-i}; t_i) \geq u_i(a_i, a_{-i}; t_i),
		\end{equation*}
		(i.e., $s_i^*(t_i)$ weakly dominates every $a_i \in A_i$), then $(s_1^*, \ldots, s_n^*)$ is a Bayesian Nash equilibrium.}

	% CONTENT FOR COLUMN 4 WILL GO HERE (Chapter 4, etc.)
	\heading{4. Dynamic Games of Incomplete Info}

	\subheading{Perfect Bayesian Equilibrium}

	\definition{A \textbf{Perfect Bayesian Equilibrium} consists of strategies and beliefs satisfying:
		\begin{enumerate}
			\item \textbf{Sequential Rationality}: Given their beliefs, players' strategies must be sequentially rational.
			      \begin{itemize}
				      \item At each information set, the action and subsequent strategy of the player must be optimal, given the player's belief and other players' subsequent strategies.
			      \end{itemize}
			\item \textbf{Consistency}: At every information set, beliefs are determined by Bayes' rule and players' strategies, whenever possible.
		\end{enumerate}
	}

	\textbf{Remark:} The strategy profile in any perfect Bayesian equilibrium is a Nash equilibrium.

	\subheading{Signaling Games}

	\textbf{Pooling vs. Separating:}
	\begin{itemize}[leftmargin=10pt,itemsep=0pt]
		\item \textbf{Pooling}: Each type sends the same message
		\item \textbf{Separating}: Each type sends a different message
		\item \textbf{Semi-separating/Partial pooling}
	\end{itemize}

	\definition{A \textbf{pure-strategy perfect Bayesian equilibrium} in a signaling game consisting of Sender's strategy $m^*(t_i)$, Receiver's strategy $a^*(m_j)$, Belief $\mu(t_i|m_j)$ satisfying:
		\begin{enumerate}[leftmargin=10pt, itemsep=0pt]
			\item \textbf{Sequential rationality}:
			      \begin{itemize}[leftmargin=10pt,itemsep=0pt]
				      \item For each $m_j \in M$, $a^*(m_j)$ maximizes the Receiver's expected utility:
				            \[a^*(m_j) = \arg\max_{a_k \in A} \sum_{t_i \in T} \mu(t_i|m_j)U_R(t_i, m_j, a_k)\]
				      \item For each $t_i \in T$, $m^*(t_i)$ maximizes the Sender's utility:
				            \[m^*(t_i) = \arg\max_{m_j \in M} U_S(t_i, m_j, a^*(m_j))\]
			      \end{itemize}
			\item \textbf{Consistency}: Given Sender's strategy $m^*$, for each $m_j$ with $T_j = \{t_i \in T : m^*(t_i) = m_j\} \neq \emptyset$:
			      \[\mu(t_i|m_j) = \frac{P(t_i)}{\sum_{t \in T_j}P(t)} \textrm{ for all } t_i \in T_j\]
		\end{enumerate}
	}


	% CONTENT FOR COLUMN 5 WILL GO HERE (Chapter 5, etc.)
	% CHAPTER 5: COOPERATIVE GAMES
	\heading{5. Cooperative Games}

	\subheading{Two-Person Cooperative Games}\\
	\definition{A pair $\Gamma = (H, d)$ where:
		\begin{itemize}[leftmargin=10pt,itemsep=0pt]
			\item $H \subset \mathbb{R}^2$ is compact and convex
			\item $d \in H$ is the threat point (payoffs if no agreement)
			\item $H$ contains at least one element, $u$, such that $u \gg d$
		\end{itemize}}

	\definition[Pareto optimality]{Payoff pairs which are not dominated by any other pair are said to be Pareto optimal. Pair $(u, v)$ dominates $(u', v')$ if $u \geq u'$ and $v \geq v'$.}

	\definition[Nash bargaining solution]{A mapping $f : W \rightarrow \mathbb{R}^2$ that associates a unique element $f(H, d) = (f_1(H, d), f_2(H, d))$ with the game $(H, d) \in W$, satisfying:
		\begin{enumerate}[leftmargin=10pt,itemsep=0pt]
			\item \textbf{Feasibility}: $f(H, d) \in H$
			\item \textbf{Individual rationality}: $f(H, d) \geq d$
			\item \textbf{Pareto optimality}: $f(H, d)$ is Pareto optimal
			\item \textbf{Invariance under linear transformations}: If a game is transformed by affine functions, the solution transforms accordingly
			\item \textbf{Symmetry}: If $(H, d)$ is symmetric, then $f_1(H, d) = f_2(H, d)$
			\item \textbf{Independence of irrelevant alternatives}: If $(H, d),(H', d') \in W$, $d = d'$, $H \subset H'$, and $f(H', d') \in H$, then $f(H, d) = f(H', d')$
		\end{enumerate}}

	\theorem{Theorem 1}{A game $(H, d) \in W$ has a unique Nash solution $u^* = f(H, d)$ satisfying the above conditions if and only if $(u^*_1 - d_1)(u^*_2 - d_2) > (u_1 - d_1)(u_2 - d_2)$ for all $u \in H$, $u \geq d$, and $u \neq u^*$.}

	\subheading{n-person Cooperative Games}\\
	\definition{For an $n$-person game with players $N = \{1, 2, \ldots, n\}$, the \textbf{characteristic function} $v$ gives each coalition $S \subseteq N$ the amount $v(S)$ that coalition $S$ can be sure of receiving. Game denoted as $\Gamma = (N, v)$.}

	\definition[Assumption]{The characteristic function $v$ satisfies:
		\begin{itemize}[leftmargin=10pt,itemsep=0pt]
			\item $v(\emptyset) = 0$
			\item For any disjoint coalitions $K$ and $L$ contained in $N$, $v(K \cup L) \geq v(K) + v(L)$ (super-additivity)
		\end{itemize}}

	\definition[Imputation]{A payoff vector $x = (x_1, \ldots, x_n)$ satisfying:
		\begin{itemize}[leftmargin=10pt,itemsep=0pt]
			\item $\sum_{i=1}^n x_i = v(N)$ (group rationality)
			\item $x_i \geq v(\{i\})$ for all $i \in N$ (individual rationality)
		\end{itemize}
		$I(N, v)$ denotes the set of all imputations.}

	\definition[Domination]{Let $x, y \in I(N, v)$, and $S$ be a coalition. We say $x$ dominates $y$ via $S$ (notation $x \succ_S y$) if:
		\begin{itemize}[leftmargin=10pt,itemsep=0pt]
			\item $x_i > y_i$ for all $i \in S$
			\item $\sum_{i \in S} x_i \leq v(S)$
		\end{itemize}
		We say $x$ dominates $y$ (notation $x \succ y$) if there is any coalition $S$ such that $x \succ_S y$.}

	\definition[Core]{The set of all undominated imputations for a game $(N, v)$ is called the core, denoted by $C(N, v)$.}

	\theorem{Theorem 2}{The core of a game $(N, v)$ is the set of all $n$-vectors, $x$, satisfying:
		\begin{itemize}[leftmargin=10pt,itemsep=0pt]
			\item $\sum_{i \in S} x_i \geq v(S)$ for all $\emptyset \neq S \subset N$
			\item $\sum_{i \in N} x_i = v(N)$
		\end{itemize}}

	\subheading{Special Games}\\
	\definition[Constant-sum game]{A game $(N, v)$ is constant-sum if for all $S \subset N$, $v(S) + v(N - S) = v(N)$.}

	\definition[Essential game]{A game $(N, v)$ is essential if $v(N) > \sum_{i \in N} v(\{i\})$. It is inessential otherwise.}

	\theorem{Theorem 3}{If $(N, v)$ is an inessential game, then for any coalition $S$, $v(S) = \sum_{i \in S} v(\{i\})$.}

	\theorem{Theorem 4}{If a game $(N, v)$ is constant-sum, then its core is either empty or a singleton $\{v(\{1\}), \ldots, v(\{n\})\}$.}

	\subheading{Strategical Equivalence}\\
	\definition[Strategical equivalence]{Two games $v$ and $u$ are strategically equivalent if there exist constants $a > 0$ and $c_1, \ldots, c_n$, such that for every coalition $S$, $u(S) = av(S) + \sum_{i \in S} c_i$.}

	\definition[Reduced form]{A characteristic function $v$ is in $(0, 1)$-reduced form if $v(\{i\}) = 0$ for all $i \in N$ and $v(N) = 1$.}

	\subheading{The Shapley Value}\\
	\definition[Shapley value]{Given any $n$-person game $(N, v)$, the Shapley value is an $n$-vector, denoted by $\phi(v)$, where the $i$-th component is:
		\begin{equation*}
			\phi_i(v) = \sum_{S \subseteq N \setminus \{i\}} \frac{s!(n-s-1)!}{n!}[v(S \cup \{i\}) - v(S)]
		\end{equation*}
		where $s = |S|$.}

	\definition[Properties of Shapley value]{The Shapley value has the following properties:
		\begin{itemize}[leftmargin=10pt,itemsep=0pt]
			\item \textbf{Individual rationality}: $\phi_i(v) \geq v(\{i\})$ for every $i \in N$
			\item \textbf{Efficiency}: $\sum_{i \in N} \phi_i(v) = v(N)$
			\item \textbf{Symmetry}: If $v(S \cup \{i\}) = v(S \cup \{j\})$ for every coalition $S$ not containing $i$ and $j$, then $\phi_i(v) = \phi_j(v)$
			\item \textbf{Additivity}: If $v$ and $w$ are characteristic functions, then $\phi(v + w) = \phi(v) + \phi(w)$
			\item \textbf{Dummy}: If a player $i$ is such that $v(S \cup \{i\}) = v(S)$ for every $S$ not containing $i$, then $\phi_i(v) = 0$
		\end{itemize}}
\end{multicols}


\end{document}
